% =====================================================================
% FASE 8 — Texto para pegar en informe_APA7.tex (Overleaf).
% DÓNDE: justo ANTES de la línea "% ================= RESULTADOS".
% Imágenes: subir a la carpeta imagenes/ de Overleaf las de docs/evidencias/ que empiezan con "fase8_"
% (fase8_comprobante_compartido.png, fase8_pago_por_cedula.png, fase8_admin_tiempo_real.png y
% fase8_pago_gas_cedula.png). Las tres primeras ya estaban subidas.
% Además (opcional): agregar en la tabla de fases la fila
%     8 & Comprobantes, pago por cédula y administración de servicios & 24/10/2026 & Completada\\
% y en la tabla de commits la fila del commit que hagas con estos cambios (código y fecha).
% =====================================================================

% ================= FASE 8 =================
\subsection{Fase 8: Comprobantes, Pago por Cédula y Administración de Servicios}

La última fase acercó la aplicación a la experiencia de una billetera comercial con tres funciones: comprobantes que se pueden compartir, la consulta de facturas por número de cédula y un panel en el que el usuario laboratorio actúa como la empresa de servicios básicos. La Tabla~\ref{tab:fase8} resume los cambios en el servidor y en la aplicación.

\begin{table}[H]
  \centering
  \caption{Funciones Añadidas en la Fase 8}\label{tab:fase8}
  \singlespacing
  \begin{tabularx}{\textwidth}{@{}>{\raggedright\arraybackslash}p{1.35in} R R@{}}
    \toprule
    \textbf{Función} & \textbf{Servidor} & \textbf{Aplicación}\\
    \midrule
    Comprobantes & Sin cambios: el N.º de comprobante es el identificador de la transacción en MySQL & Panel modal de detalle en cada movimiento con el botón Ver comprobante, y comprobante automático al terminar una operación; se comparte como imagen PNG mediante FileProvider\\
    Pago por cédula & Columnas cedula y periodo en la tabla contratos (migración); GET /facturas?cedula=\&servicio= con validación de la cédula & Modo Por cédula en el diálogo de pago: lista de facturas pendientes y pagadas; se elige una y se paga\\
    Administración de servicios & Tabla servicios (catálogo); GET /servicios; GET y POST /admin/facturas y POST /admin/servicios solo para el laboratorio (403 para otras cuentas); aviso FACTURAS\_ACTUALIZADAS por WebSocket & Pantalla de inicio propia del administrador; panel con pestañas Facturas y Servicios: resumen, filtros, emisión de facturas, alta de servicios y actualización en tiempo real\\
    \bottomrule
  \end{tabularx}
  \nota{Elaboración propia.}
\end{table}

\subsubsection{Comprobantes Compartibles}

Al tocar un movimiento del historial se abre un panel modal inferior con la descripción de la operación (detalle, categoría, fecha y hora, número de comprobante y estado) y un botón Ver comprobante; así la lista se mantiene compacta aunque haya muchos movimientos. Ese botón abre un comprobante con el monto, la fecha y hora, el número de comprobante, la cuenta, el titular y el detalle de la operación. Tras una operación exitosa el comprobante se abre automáticamente, como en las aplicaciones bancarias. Como sus datos provienen de Room, el comprobante se puede consultar y compartir sin conexión. Para compartirlo, la interfaz se dibuja dentro de una capa gráfica de Compose que luego se convierte en imagen; el archivo PNG se guarda en la caché privada de la aplicación en Dispatchers.IO y se entrega a la aplicación elegida por el usuario mediante FileProvider, que otorga un permiso de lectura temporal solo a esa aplicación (Figura~\ref{fig:comprobante}). El botón Descargar guarda la misma imagen en la Galería del teléfono, en el álbum Billetera CC: desde Android 10 se usa MediaStore, que no requiere permisos, y en Android 8 y 9 se solicita el permiso de almacenamiento antes de escribir en la carpeta pública de imágenes.

\begin{figure}[H]
  \centering
  \caption{Comprobante Compartido como Imagen}\label{fig:comprobante}
  \includegraphics[width=0.45\textwidth,keepaspectratio]{imagenes/fase8_comprobante_compartido.png}
  \nota{Imagen PNG generada por la aplicación y enviada al menú de compartir de Android.}
\end{figure}

\subsubsection{Pago de Servicios por Cédula}

Las facturas de servicios básicos ahora registran la cédula del titular y el mes facturado. En el diálogo de pago, además del código de contrato, el usuario puede escribir una cédula, elegir el servicio y ver todas las facturas de esa persona: las pendientes primero y, en las pagadas, la fecha y la cuenta que realizó el pago (Figura~\ref{fig:pagocedula}). La cédula se valida en la aplicación y de nuevo en el servidor con el algoritmo de módulo 10 de la cédula ecuatoriana: provincia válida, tercer dígito menor que seis y dígito verificador calculado con los coeficientes 2, 1, 2, 1, 2, 1, 2, 1, 2. Las cédulas de los datos de demostración son ficticias, aunque cumplen el algoritmo, y el validador se cubrió con pruebas unitarias.

\begin{figure}[H]
  \centering
  \caption{Búsqueda de Facturas por Cédula en el Diálogo de Pago}\label{fig:pagocedula}
  \includegraphics[height=0.5\textheight,keepaspectratio]{imagenes/fase8_pago_por_cedula.png}
  \nota{La factura elegida se resalta y el botón muestra el monto a pagar.}
\end{figure}

\subsubsection{Administración de Servicios en Tiempo Real}

El usuario laboratorio no es un cliente sino el administrador del sistema: representa a la empresa de servicios básicos y al laboratorio de pruebas. Por ello tiene una pantalla de inicio propia, sin saldo ni botones de operación, que muestra el estado del túnel WebSocket, los nodos conectados, el total por cobrar y el total recaudado, y dos accesos: Servicios y facturas, y Pruebas de concurrencia. El panel de servicios tiene dos pestañas. En Facturas, la empresa emite cada mes la factura de una persona indicando su cédula, el servicio, el mes y el valor; el servidor genera el código de contrato y esa persona la encuentra después al buscar su cédula en el pago de servicios. En Servicios se administra el catálogo, que pasó de estar fijo en el código a una tabla de MySQL: el laboratorio puede agregar servicios nuevos, como gas o telefonía, con un código y un ícono, y estos aparecen de inmediato al emitir facturas y en el diálogo de pago de los clientes (Figura~\ref{fig:pagogas}). El servidor rechaza con el código 403 cualquier intento de administración desde otra cuenta. Cuando un cliente paga una factura desde cualquier dispositivo, el servidor envía al laboratorio el aviso FACTURAS\_ACTUALIZADAS por WebSocket; el panel recarga la lista sin intervención del usuario y muestra quién pagó (Figura~\ref{fig:admintiemporeal}). Este comportamiento extiende al dominio de los servicios el mismo mecanismo de notificación en tiempo real usado para las transferencias.

\begin{figure}[H]
  \centering
  \caption{Panel de Administración Actualizado en Tiempo Real}\label{fig:admintiemporeal}
  \includegraphics[height=0.5\textheight,keepaspectratio]{imagenes/fase8_admin_tiempo_real.png}
  \nota{El aviso inferior y la factura marcada como pagada aparecieron sin recargar, cuando la cuenta 1002 pagó desde otro cliente.}
\end{figure}

\begin{figure}[H]
  \centering
  \caption{Factura de un Servicio Nuevo Encontrada por el Cliente}\label{fig:pagogas}
  \includegraphics[height=0.5\textheight,keepaspectratio]{imagenes/fase8_pago_gas_cedula.png}
  \nota{El laboratorio agregó el servicio GAS y emitió una factura; el cliente la encuentra al buscar su cédula en el diálogo de pago.}
\end{figure}

